\documentclass[11pt]{article}
\usepackage{fontspec}
\usepackage[french,english]{babel}
\usepackage{amsmath}
\usepackage{amssymb}
\usepackage{amsthm}
\usepackage{mathrsfs}
\usepackage{tikz-cd}
\usepackage{enumitem}

\newcommand{\printedpage}[1]{\marginpar{\footnotesize p.~#1}}
\newcommand{\msp}[1]{\marginpar{\footnotesize [#1]}}
% printed notation: the category C' enclosed in a V-shaped wedge with the arrow name under it
\newcommand{\Vcat}[2]{\underset{#2}{\underset{\vee}{#1}}}

\begin{document}

% ===================== PDF page 1 =====================
\printedpage{--}
% no page number printed (title page)
% head: (none)

\begin{center}
\vspace*{4cm}
{\Large\textsf{THE DERIVATORS}}

\bigskip\bigskip

\textbf{ALEXANDRE GROTHENDIECK}

\bigskip\bigskip\bigskip

\textbf{Chapter II}

\textbf{Cofinality (on the right and on the left)}

\textbf{(preliminary to cohomological cofinality)}
\end{center}

\clearpage

% ===================== PDF page 2 =====================
\printedpage{--}
% no page number printed (editorial note)
% head: (none)

\vspace*{6cm}

\noindent This text was deciphered and transcribed in \LaTeXe{} by M. Künzer. It was edited by M. Künzer, J. Malgoire, and G. Maltsiniotis. The transcription is as faithful as possible to the manuscript. For the few obvious corrections, or rare comments by the editors, as well as for the original numbering of the manuscript pages, typewriter type \texttt{[typewriter]} in brackets is used. A question mark in brackets means that one is not sure of the preceding word.

\noindent This edition is provisional. Remarks, comments and corrections are welcome. Send a message to:

\noindent\textbf{[e-mail address removed]}

\hfill G. Maltsiniotis\hspace*{2cm}

\clearpage

% ===================== PDF page 3 =====================
\printedpage{1}
% head: A. Grothendieck : Dérivateurs, Ch. II, Édition des Universités de Montpellier II et Paris VII
% foot: Version 8 juillet 2001

\msp{page 1}

\section*{1\quad Cohomological cofinality}

Let
\[ \alpha : C' \longrightarrow C \]
be a functor of (small) categories, $\mathcal{M}$ a category having the relevant $\varprojlim$, $F$ a presheaf (thus contravariant in $i$) on $C$ with values in $\mathcal{M}$. We have a canonical homomorphism in $\mathcal{M}$
\[ \varprojlim_{C} F \longrightarrow \varprojlim_{C'} \underbrace{\alpha^*(F)}_{F\circ\alpha} \tag*{(1)} \]
and we want conditions on $\alpha$ for this always to be an isomorphism (\footnote{One then says that $\alpha$ is \textit{cofinal,} cf. SGA 4, I 8.1, or 3.3, p. 45.}). We are at once reduced to the case where $\mathcal{M} = (\mathrm{Ens})$. Let us say that a category $I$ is \textit{weakly left filtering} if $I \neq \emptyset$ and if $\forall\, i', i'' \in \operatorname{Ob} I$, $\exists\, i \in \operatorname{Ob} I$ and arrows
\[
\begin{tikzcd}[row sep=small, column sep=small]
 & i' \\
i \arrow[ur] \arrow[dr] & \\
 & i''.
\end{tikzcd}
\]
This also means that $\operatorname{Ob} I$ as an ordered set is left filtering.

\textbf{Proposition.} \textit{For} (1) \textit{always to be an isomorphism} ($\alpha$ \textit{``anodyne for} $\varprojlim$ \textit{''), it suffices that for every} $x_0 \in \operatorname{Ob} C$\textit{, the category} $x_0\backslash C'$ \textit{of pairs} $(x', u)$\textit{, with} $x' \in \operatorname{Ob} C'$\textit{,} $u : \alpha(x') \longrightarrow x$\textit{, be weakly left filtering} (\footnote{Cf. improvement in the remarks p. 5.}, \footnote{Cf. 3.2, p. 43-44, for necessary and sufficient conditions for cofinality.}). % [sic] u : \alpha(x') \longrightarrow x -> x_0 in place of x, the category is x_0\backslash C'

This does not imply the stronger conclusion
\[ F \underbrace{\xrightarrow{\;\sim\;}}_{\text{adj.}} \alpha_*\alpha^* F \]
that the canonical adjunction morphism is an isomorphism. For this to be so, it suffices that $\forall\, x_0 \in C$, the functor $\alpha/x_0 : C'/x_0 \longrightarrow C/x_0$ induced by $\alpha$ be ``anodyne for the $\varprojlim$ '', which also means that for every arrow $u : x_1 \longrightarrow x_0$ in $C$, the category $\Vcat{C'}{u}$ formed of triples $(x', \xi, \eta)$ with $x' \in \operatorname{Ob} C'$ and $\xi : x_1 \longrightarrow \alpha(x')$, $\eta : \alpha(x') \longrightarrow x_0$ rendering commutative
% printed: C' inside a V-shaped wedge, with u printed below the wedge
\msp{page 2}
the diagram
\printedpage{2}
% head: A. Grothendieck : Dérivateurs, Ch. II, Édition des Universités de Montpellier II et Paris VII
% foot: Version 8 juillet 2001
\[
\begin{tikzcd}
x_1 \arrow[rr, "u"] \arrow[dr, "\xi"'] & & x_0 \\
 & \alpha(x'), \arrow[ur, "\eta"'] &
\end{tikzcd}
\]
is weakly left filtering.

More generally, let us consider
\[
\begin{tikzcd}
X' = C' \arrow[rr, "\alpha"] \arrow[dr, "f'"'] & & C = X \arrow[dl, "f"] \\
 & S &
\end{tikzcd}
\]
and let us ask ourselves the question when the functorial homomorphism, with respect to a variable presheaf $F$ on $X$,
\[ f_*(F) \longrightarrow f'_*(\alpha^*(F)) \]
is always an isomorphism. (The preceding case is the one where $S = X$, $f = \mathrm{id}_X$.) We find that it suffices that $\forall\, s_0 \in \operatorname{Ob} S$, the functor
\[ \alpha/s_0 : X'/s_0 \longrightarrow X/s_0 \]
be ``anodyne for the $\varprojlim$ '', and for this, it suffices that for \texttt{[every]} object $(x_0, u_0)$ of $X/s_0$ (where
\[ f(x_0) \xrightarrow{u_0} s_0) \]
the category $(x_0, u_0)\backslash(X'/s_0)$ be weakly left filtering.

\msp{page 3}

This category is interpreted as that of triples $(x', \xi, u)$ with
\[
\left\{
\begin{aligned}
& x' \in \operatorname{Ob} C' \\
& u : g(x') = f(\alpha(x')) \longrightarrow s_0 \\
& \xi : x_0 \longrightarrow \alpha(x')
\end{aligned}
\right.
\]
\texttt{[rather} $f'(x')$ \texttt{instead of} $g(x')$\texttt{]} so as to render commutative the diagram
\[
\begin{tikzcd}
f(x_0) \arrow[rr, "u_0"] \arrow[dr, "f(\xi)"'] & & s_0 \\
 & f(\alpha(x')). \arrow[ur, "u"'] &
\end{tikzcd}
\]
I would like to bring out variants of these statements in the cohomological setting, the functors $\varprojlim$ being replaced by $R\Gamma$'s (of abelian presheaves e.g. - but one also thinks of the ``noncommutative cases'', defined by derivators in particular.)

\clearpage

% ===================== PDF page 5 =====================
\printedpage{3}
% head: A. Grothendieck : Dérivateurs, Ch. II, Édition des Universités de Montpellier II et Paris VII
% foot: Version 8 juillet 2001

If $I$ is a \textit{finite ordered} category, the \texttt{[weakly]} left filtering condition is equivalent to the existence of an initial object, which implies that $I$ is contractible. In the preceding statements, it might be that the right philosophy is to replace ``weakly left filtering'' everywhere by ``contractible''. If e.g. $C$ is the final category, then replace the condition that $C'$ be left filtering, by acyclicity. This is in any case in accord with the spirit of the axiom $\mathrm{Der}_5$ (case where $J$ is the final category). Indeed, in the case where $C = e$, the condition that $\alpha$ be anodyne for the $\varprojlim$ is visibly equivalent to saying that the functor $\alpha^* : \mathcal{M} \longrightarrow \underline{\mathrm{Hom}}(C'^{o}, \mathcal{M})$, associating to every $F \in \mathcal{M}$ the constant presheaf with value $F$ % [sic] précedents -> précédents, missing accent on the second e
\msp{page 4}
on $C'$, gives an adjunction morphism
\[ \varprojlim_{e} F = F \longrightarrow \alpha_*(\alpha^*(F)) = \varprojlim_{C'} \alpha^* F \]
which is an isomorphism. And the axiom $\mathrm{Der}_5$ ensures that it must be so if $C'$ is contractible.

The filtering condition in the proposition is moreover a \textit{sufficient} condition for anodyneness, but by no means necessary. If e.g. $C$ is the final category, then the condition of anodyneness for $C' \longrightarrow C$ is simply that $C'$ be connected and nonempty, i.e. $\pi_0(C') = \{0\}$. (Condition which visibly ignores the higher $\pi_i$.)

Let \fbox{$\alpha : X' \longrightarrow X$} be a functor of small categories. The condition suggested by the proposition, that the $x_0\backslash X'$ be contractible, means that $\alpha$ is a co-strict hot-equivalence, i.e. that $\alpha^{o} : X'^{o} \longrightarrow X^{o}$ is a hot-equivalence ``locally on $X^{o}$'', hence also globally. If $\alpha$
\msp{page 5}
itself were a strict hot-equivalence, instead of $\alpha^{o}$, then we would have by $\mathrm{Der}_5$, that for $F \in \mathbf{D}(X)$, $F \xrightarrow{\;\sim\;} \alpha_*\alpha^* F$ ($\mathbf{D}$ a left derivator), which is much stronger than $\Gamma_{\mathbf{D}}(X, F) \xrightarrow[\sim]{\alpha^*} \Gamma_{\mathbf{D}}(X', \alpha^* F)$. The dual condition implies (if $\mathbf{D}$ \texttt{[is a]} right derivator) $\alpha_!\alpha^* F \xrightarrow{\;\sim\;} F$. But is it reasonable to expect that it also implies the condition ``on the cohomology side, instead of homology'',
\[ \Gamma_{\mathbf{D}}(X, F) \xrightarrow{\;\sim\;} \Gamma_{\mathbf{D}}(X', \alpha^*(F)) \ \ ? \]

The first case, which should carry conviction, is that of the derivator DerAb, i.e. of the derived categories of abelian sheaves:

\textbf{Question.} \textit{Let} $\alpha : X' \longrightarrow X$ \textit{be a co-strict} hot-\textit{equivalence,} $F$ \textit{an abelian presheaf on} $X$\textit{. Do we then have} $H^i(X, F) \xrightarrow[\sim]{\alpha^*} H^i(X', \alpha^*(F))$ \textit{for every} $i \in \mathbf{Z}$ \textit{?}

If $X$ admits an initial object $\emptyset$, then $\emptyset\backslash X' \simeq X'$, so $X'$ contractible, so it's OK for $F$ constant.

\textbf{Remark.} I realize that in proposition 1, instead of the condition that the $x\backslash C'$ be weakly left filtering, it would suffice that they be nonempty and connected, i.e.

\end{document}
